
\documentclass[tikz,border=2pt]{standalone}
\usepackage[utf8]{inputenc}
\usepackage[T1]{fontenc}
\usepackage[french]{babel}
\usepackage{amsmath,amssymb}
\usepackage[table]{xcolor}
\usepackage{siunitx}
\usepackage[version=4]{mhchem}
\usepackage{tabularx,booktabs,array}
\usepackage{tikz}
\usetikzlibrary{arrows.meta,positioning,calc,fit,decorations.pathmorphing,patterns,shapes.geometric}
\usepackage{pgfplots}
\pgfplotsset{compat=1.18}
\definecolor{primary}{RGB}{14,116,144}
\definecolor{accent}{RGB}{16,185,129}
\definecolor{dark}{RGB}{15,23,42}
\definecolor{bglight}{RGB}{241,245,249}
\definecolor{borderlight}{RGB}{203,213,225}
\definecolor{examplepink}{RGB}{236,72,153}
\definecolor{remarkblue}{RGB}{14,165,233}
\definecolor{notepurple}{RGB}{99,102,241}
\definecolor{synthgreen}{RGB}{5,150,105}
\definecolor{synthorange}{RGB}{249,115,22}
\definecolor{synthviolet}{RGB}{79,70,229}
\definecolor{primary}{RGB}{14,116,144}
\definecolor{accent}{RGB}{16,185,129}
\definecolor{dark}{RGB}{15,23,42}
\definecolor{bglight}{RGB}{241,245,249}
\definecolor{examplepink}{RGB}{236,72,153}
\definecolor{remarkblue}{RGB}{14,165,233}
\definecolor{notepurple}{RGB}{99,102,241}
\definecolor{synthgreen}{RGB}{5,150,105}
\definecolor{synthorange}{RGB}{249,115,22}
\definecolor{synthviolet}{RGB}{79,70,229}
\begin{document}
\begin{tikzpicture}[
 node distance=0.8cm and 1.2cm,
 every node/.style={font=\small},
 box/.style={rectangle, rounded corners=3pt, draw=#1, fill=#1!8, thick,
 text width=5.5cm, minimum height=1.2cm, align=center, font=\small},
 arrow/.style={-{Stealth[length=3mm]}, thick, #1}
]

% Titre
\node[font=\large\bfseries, color=dark] (titre) at (0,5.5) {Sch\'{e}ma bilan -- Mod\`{e}le du gaz parfait};

% Boîte centrale
\node[box=primary, text width=6cm, minimum height=1.6cm, font=\normalsize\bfseries] (eq) at (0,3.5)
 {$PV = nRT$ \\ \'{E}quation d'\'{e}tat du gaz parfait};

% Hypothèses (gauche)
\node[box=synthorange, text width=5.2cm] (hyp) at (-5.5,3.5)
 {\textbf{Hypoth\`{e}ses} \\ Entit\'{e}s ponctuelles \\ Pas d'interaction \`{a} distance};

% Grandeurs (droite)
\node[box=synthviolet, text width=5.2cm] (grand) at (5.5,3.5)
 {\textbf{Grandeurs} \\ $P$ (Pa), $V$ (m$^3$), $n$ (mol), \\ $T$ (K), $R = \SI{8,314}{\joule\per\mol\per\kelvin}$};

% Volume molaire (bas gauche)
\node[box=synthgreen, text width=5.2cm] (vm) at (-4,0.5)
 {\textbf{Volume molaire} \\ $V_m = \dfrac{RT}{P}$ \\ Ind\'{e}pendant de la nature du gaz};

% Avogadro (bas droite)
\node[box=examplepink, text width=5.2cm] (avo) at (4,0.5)
 {\textbf{Loi d'Avogadro-Amp\`{e}re} \\ M\^{e}mes $P$, $V$, $T$ \\ $\Rightarrow$ m\^{e}me $n$};

% Limites (bas centre)
\node[box=red!70!black, text width=6cm] (lim) at (0,-1.5)
 {\textbf{Limites} \\ $P < \SI{1}{\mega\pascal}$ \quad et \quad $T$ pas trop basse};

% Flèches
\draw[arrow=synthorange] (hyp) -- (eq);
\draw[arrow=synthviolet] (eq) -- (grand);
\draw[arrow=synthgreen] (eq.south west) -- +(-0.5,-0.5) -- (vm.north);
\draw[arrow=examplepink] (eq.south east) -- +(0.5,-0.5) -- (avo.north);
\draw[arrow=red!70!black] (eq.south) -- (lim.north);

\end{tikzpicture}
\end{document}
